\scanpage{69}
\mspage{68}
% E: text|top left|A Recollements donnant lieu
A) Recollements donnant lieu à une application naturelle
% E: display|upper left|Xcal to X
\[\mathcal X\longrightarrow X,\]
% E: text|left upper middle|ou ce qui revient au même
ou ce qui revient au $=$, où l'isom. $\sigma:\alpha\longrightarrow\beta$ \inserted{(\unsure{moyennant au préalable|moyennant éventuellement} l'antipodisme)} est obtenu par transport parallèle --- ou encore, \unsure{combinatoirement}, \unsure{car où|où} l'isom. de transition $\omega_D\mathbin{-}\omega_{D'}$ est opposée au transport parallèle. Il y a deux possibilités (avec tjrs $\beta=b$)
% E: text|left middle|1 Cas alpha=a'
1) Cas $\alpha=a'$ --- on recolle \deleted{l'arc de \unsure{la p.r.}} \deleted{une rive \unsure{$b$ dans}} \inserted{l'arc de spéc $b$ sur $\Delta_D$} sur l'arc \deleted{$a'$ qui est} de gén. $a'$ sur $\Delta_{D'}$ qui est sur la rive \uline{faisant face} à celle de $b$, i.e. opposée à celle obtenue par transport parallèle: \deleted{C'est} Alors $\mathcal X\longrightarrow X$ est une immersion. \deleted{On repère}
% E: margin note|left margin at Pour définir|emphasis mark
\margin{\unsure{*}}
% E: text|left lower middle|Pour définir la structure polygonale
Pour définir la structure polygonale sur l'arc de transition à partir d'une p.r. donnée $D'$, on repère \deleted{\unsure{d'une op.}} une générisation par la rive qui est \deleted{\unsure{dans la direc.}} celle dans laquelle on va déplacer $D'$ pour génériser: le trait bleu qui marque la rive ``tire'' sur $D'$ pour le déplacer vers lui.
% E: text|bottom left continuing right|2 Cas alpha=a
2) Cas $\alpha=a$, --- on recolle $b$ sur l'arc $a$ que l'on déduit par transport parallèle. Maintenant $\mathcal X\longrightarrow X$ est une immersion sur chaque face $F^*$, mais c'est un rebroussement le long d'une frontière commune à deux faces.

On repère une op. de générisation par \deleted{la rive l'} l'arc de gén. qui est opposé à la rive dans laquelle on se \unsure{pousse} --- le trait bleu marquant \unsure{cette rive va venir} ``pousser'' sur $D'$ pour le génériser, et le \unsure{déloger|dilater} d'un sommet envisagé.
